% 9_conclusion.tex — 结论

本文围绕烟幕干扰弹的投放策略优化问题，建立了完整的数学模型并进行了系统求解，主要工作和结论如下：

\textbf{（1）建立了完整的运动学和几何模型。}从导弹、无人机、干扰弹和烟幕云团的运动学出发，构建了基于视线遮挡几何判定的烟幕遮蔽效果评估模型，为投放策略的定量优化提供了理论基础。

\textbf{（2）设计了分层递进的优化求解框架。}针对5个复杂度递增的子问题，分别采用直接仿真（问题1）、网格搜索+SQP混合优化（问题2）和遗传算法+SQP精化（问题3--5）的求解策略，在保证求解质量的同时控制了计算复杂度。

\textbf{（3）揭示了关键参数的影响规律。}通过灵敏度分析发现，烟幕半径是影响遮蔽效果最敏感的参数，无人机速度约束限制了部署灵活性，导弹速度变化体现了威胁对抗的动态特性。

\textbf{（4）给出了具有工程参考价值的最优投放策略。}各问题的优化结果以表格形式输出，可直接用于指导无人机的任务规划。问题5的预分配策略为大规模多UAV-多导弹场景提供了一种可扩展的求解思路。

\textbf{（5）展望了模型改进方向。}指出引入随机烟幕扩散模型、连续遮蔽度函数和动态重规划能力是提升模型实战适用性的关键方向。

综上，本文模型结构完整、求解方法得当、结果分析详尽，对烟幕干扰弹的战术运用具有一定的理论指导意义和工程参考价值。
